\section{模型验证、评价与适用边界}\label{sec:validation}
本文的可信度来自可追溯的计算证据，而非算法名称。验证依次回答三个问题：计算是否针对同一套输入；输出是否真的满足规定约束；通过检查后还能声称多强的结论。四问的结果口径先汇总于表\ref{tab:summary}，其中累计作业与并行完工必须分别理解。
\papertable{summary}{四问使用同一物理口径；Q1 为累计作业，Q2 为运输完工，Q3 为联合完工，Q4 资源配置另见表\ref{tab:resources}}{summary}

\subsection{从数据版本到逐箱逐资源回放}
输入版本一致性核验防止将旧输出与新参数混用，也防止 Q4 从不同输入版本的 Q3 结果继续分区。论文诊断另记录正式输出摘要，生成宏、表和图时检查来源是否变化。这保证证据来自当前版本，但文件摘要本身不证明物理公式正确。

随后不直接相信结果表中的能耗、SOC 和交付时刻，而从官方箱单、访问顺序与机型重算：每箱恰好一次，卸货后更新载荷，空载返航，逐航段累计能量，交接结束才视为交付。资源核验分别检查无人机返航释放、电池充满释放、中继周转与组件充电；时限核验区分硬截止和软延误。Q2 保存的全部候选也逐一回放，Q4 则重新枚举固定计划的分区。

为检验验证器能否拒绝错误，使用输入篡改、输出能耗或时刻篡改、漏箱/重箱、嵌套占用区间与容差冲突等反例。物理边界涵盖零距离、零载荷、额定载荷、恰好达到储备、DEM 边界与角点、无效像元及双向链路覆盖内外。DEM 穿越另与独立矩形裁剪算法交叉核对，避免仅用同一几何函数产生测试预期。

这些证据有明确分工：回放检查约束和表格真实性，反例检查错误是否会被发现，独立几何参照与量纲核对补充共享内核的盲区。能耗回放仍调用官方物理引擎，所以不能排除所有共享公式错误，更不能代替实际飞行校准。具体测试命令和完整记录随源码保存，正文只保留支持模型可信范围的证据。

\subsection{加密通信回放与最紧约束}
\paperfig{fig12_validation}{基准可行并不等于余量充足：最低返航 SOC 为 \MinSoc\%，最小硬截止余量为 \HardSlack{} s，最小通信余量为 \MarginHalf{} dB。三类证据分别对应能源、及时性和通信风险。}{validation}
\papertable{sampling}{修复后的同一联合计划在1 s与0.5 s回放中均无未覆盖状态}{sampling}
图\ref{fig:validation}和表\ref{tab:sampling}给出最紧约束。\HardBoxes{} 个硬截止货箱均未违反截止，其中最紧余量为\HardSlack{} s。通信 1 s 设置检查 \StatesOne{} 条记录，0.5 s 设置检查 \StatesHalf{} 条，两者均无未覆盖状态，最小余量分别为 \MarginOne{} 与 \MarginHalf{} dB。另以未参与搜索的0.25 s网格回放 \StatesQuarter{} 条状态，未覆盖数为 \ViolationsQuarter{}；该检查通过仍不代表连续时间证明。

两种网格都按各相位长度分别等分，并不严格嵌套；状态记录还包含相位边界重复点，所以既不能称为独立统计样本，也不能据两点画出一般性的收敛规律。遮挡判定可能在采样点之间突变。连续通信安全仍需要定位遮挡变化事件，或为整段时间建立可验证的余量下界；当前证据尚未完成这一步。因此，Q3 面对题设连续通信要求仍存在明确证明缺口。

\subsection{四类结论的精确性范围}
四问所称“解”的含义不同，表\ref{tab:scope}集中列出其边界。特别是，精确求解一个受限问题，不等于精确求解包含所有离散和连续选择的原问题。
\begin{table}[htbp]\centering\caption{算法证据与结论强度一一对应}\label{tab:scope}
\begin{tabularx}{\linewidth}{lXX}\toprule
问题 & 已成立的结论 & 不由当前结果证明的结论\\\midrule
Q1 & 单点往返、当前物理规则与词典序目标下，集合划分 DP 精确 & 加入实体资源、时限及多点访问后的调度最优\\
Q2 & 架次列与区间着色生成全准时候选；有限目录和网格上提供修正对偶下界 & 全局最优、完整 Pareto 前沿、所选架次数唯一最优\\
Q3 & 外层离散搜索与条件时序LP；1 s、0.5 s及独立0.25 s回放通过 & 路线、货箱、中继位置/高度的联合全局最优及连续通信安全\\
Q4 & 固定 Q3 时序及中继复制解释下，分区枚举精确 & 允许全面重规划时的全局最少资源或经济成本最优\\\bottomrule
\end{tabularx}\end{table}
\FloatBarrier

\subsection{模型优势、局限与改进顺序}
公共物理与上层决策分离，使每个结果均可沿货箱、航段、资源占用和通信状态追溯。Q1与固定Q3后的Q4可精确求解；Q2和Q3采用受限列模型、明确热启动与固定种子邻域搜索，没有原问题全局最优或统计显著性证据。官方等效能耗未经实飞辨识，也不支持天气外推。

已发现的通信反例经中继修复与重新搜索消除。后续应优先定位采样间的遮挡变化；再评估能耗误差、充电变化与最紧截止余量下的可执行性，最后扩大分组、中继与资源顺序搜索。现有候选和反例均可复算，不能替代实飞验证。
\FloatBarrier
